%==============================================================================
%  OUTLINE: Minimum-Fuel Low-Thrust GTO -> Cislunar Transfers (tulip & ELFO)
%
%  Skeleton only. Each section gives the intended content in a few sentences
%  plus bullet points. Source of record for all numbers:
%    optimal_control/orbit_transfer/GTO_tulip/LOW_THRUST_MINFUEL_CAMPAIGN.md
%    optimal_control/orbit_transfer/GTO_tulip/ROADMAP.md
%    optimal_control/orbit_transfer/GTO_tulip/sundman_minfuel_solution_note.tex
%==============================================================================
\documentclass[11pt]{article}
\usepackage[margin=1in]{geometry}
\usepackage{amsmath,amssymb}
\usepackage{enumitem}
\usepackage[capitalise]{cleveref}
\setlist{nosep}

\title{Minimum-Fuel Low-Thrust Transfers from GTO to\\
       Cislunar Tulip and ELFO Orbits\\[4pt]
       \large (Paper Outline)}
\author{M.\ Casey \and D.\ Koblick}
\date{\today}

\begin{document}
\maketitle

\begin{abstract}
\emph{[Outline note.]} One paragraph: pose the minimum-fuel low-thrust
GTO$\rightarrow$cislunar transfer, state that we certify a sharp bang-bang
solution on the full $\sim$40-revolution spiral to machine tolerance, and that
we map the $\Delta V$--transfer-time front for two operationally relevant
targets---a south-pole \emph{tulip} orbit and a lunar \emph{ELFO}. Headline
numbers: tulip $\Delta V = 3.37$~km/s (25 switches, defect $2.4\times10^{-14}$)
at $1.15\times$ min-time; ELFO front minimum $2.69$~km/s at $1.73\times$
min-time. \emph{[Outline note: the ``$\times$~min-time'' factors quoted here are
the legacy tulip-$6.2907$~ND scaling; final copy must state them against each
target's own certified min-time (tulip $\sim$5.8~ND, ELFO $6.0962$~ND) --- see
\S2/\S4. Physical $t_f$/$\Delta V$ are unaffected.]}
\end{abstract}

%==============================================================================
\section{Introduction}
%==============================================================================
\emph{[A few sentences.]} Motivate low-thrust cislunar delivery: high-$I_{sp}$
electric propulsion trades trip time for propellant, which dominates rideshare
economics from a GTO drop-off. The many-revolution min-fuel problem is the hard
member of the min-time/min-energy/min-fuel trio because its cost is \emph{linear}
in the control, so the optimum is genuinely bang-bang (burn near perigee, coast
near apogee) with dozens of switches, and it lives on a $\sim$40-revolution
CR3BP spiral whose perigee passes give $\sim10^6$ shooting sensitivity. We show
both walls can be knocked down together and the resulting solutions certified.

\begin{itemize}
  \item Problem framing: GTO $\rightarrow$ cislunar via continuous low thrust in
        the Earth--Moon CR3BP; why min-fuel matters commercially.
  \item Why min-fuel is hard: linear-in-control $\Rightarrow$ bang-bang
        (objective-side wall) compounded by the many-perigee spiral
        (dynamics-side wall). The two-wall framing is the paper's spine.
  \item Contributions: (i) machine-tight certified full-spiral bang-bang
        solution; (ii) energy$\rightarrow$fuel homotopy + Sundman regularization
        recipe; (iii) PMP cross-check of a direct solution via KKT duals;
        (iv) $\Delta V$--time fronts for two targets.
  \item Related work: Bertrand--\'Ep\'enoy continuation, Sundman/regularized
        transcriptions, prior GTO$\rightarrow$cislunar low-thrust studies.
        \emph{[cite]}
\end{itemize}

%==============================================================================
\section{Target Orbits: Tulip and ELFO}
\label{sec:constellation}
%==============================================================================
\emph{[A few sentences.]} Define the two rendezvous targets and why each is of
interest for lunar south-pole service. The \emph{tulip} is a south-pole-favorable
petal orbit (single plane, single-hop nav/relay geometry); the \emph{ELFO} is
the elliptical lunar frozen orbit used as an Intuitive-Machines-style data-relay
reference. Both are specified as fixed rendezvous states in CR3BP nondimensional
units.

\begin{itemize}
  \item \textbf{Tulip}: south-pole tulip point; role in nav/relay constellation
        (link to proj7 line of work). Give the target state used.
  \item \textbf{ELFO}: \texttt{im\_elfo\_optimum} --- sma 12{,}000~km, ecc 0.69,
        inc $56.5^\circ$, argp $90^\circ$; role as relay benchmark.
  \item Common departure: GTO $350\times35786$~km, $\omega=-25^\circ$.
  \item CR3BP setup: $\mu^\star = 0.012150585609624$; nondimensionalization;
        spacecraft 15~kg, 25~mN, $I_{sp}=2100$~s.
  \item Min-time anchors (front endpoints, 0 switches), computed as a hard
        all-burn ($s\equiv1$) free-time min-time solve on the two-primary-clock
        transcription (\S3):
        \begin{itemize}
          \item \textbf{ELFO}: $t_f=6.0962$~ND $=27.02$~d (certified: machine-tight,
                defect $1.7\times10^{-15}$, primer $0.17^\circ$).
          \item \textbf{Tulip}: $t_f=5.8267$~ND $=25.83$~d (certified:
                machine-tight, defect $1.7\times10^{-15}$, rendezvous exact,
                primer $0.19^\circ$, \emph{mesh-invariant}), to the backbone
                rendezvous state the $\Delta V$--$t_f$ front actually targets
                (dMoon $\sim$28{,}000~km, speed $0.31$~ND).
        \end{itemize}
        \emph{[Outline note --- IMPORTANT.]} The frequently quoted tulip min-time
        $t_f=6.2907$~ND ($\Delta V=4.47$~km/s) is the min-time to a \emph{different}
        tulip state (the max-$\dot y$ point, dMoon $\sim$6{,}000~km) and must
        \emph{not} be used to scale this front: the min-fuel front rendezvouses at
        the backbone state above, whose true min-time is $\sim$5.8~ND. All ``$\times$
        min-time'' factors in \S4 are therefore relative to the target-consistent
        anchors here, not to $6.2907$.
\end{itemize}

%==============================================================================
\section{The Optimal Control Problem}
%==============================================================================
Delivering a spacecraft from geostationary transfer orbit (GTO) to a cislunar
target under continuous low thrust is posed here as a fixed-time minimum-fuel
optimal control problem in the Earth--Moon circular restricted three-body
problem (CR3BP). Fixing the transfer time and minimizing propellant isolates the
quantity of engineering interest, the electric-propulsion equivalent $\Delta V$,
and turns the trade against trip time into a one-parameter sweep
(\cref{sec:results}). The problem is solved by direct collocation: the trajectory
is transcribed into a nonlinear program (NLP) and solved with an interior-point
method. This section states the continuous-time problem, identifies why its
minimum-fuel form is numerically hard, and describes the two regularizations that
make it solvable to machine tolerance.

\subsection{Continuous-time formulation}
Let $\boldsymbol r=(x,y,z)$ and $\boldsymbol v=\dot{\boldsymbol r}$ be
nondimensional (ND) position and velocity in the rotating Earth--Moon frame and
$m$ the spacecraft mass, with mass ratio $\mu=0.012150585609624$ and primaries
fixed at $\boldsymbol r_E=(-\mu,0,0)$ (Earth) and $\boldsymbol r_M=(1-\mu,0,0)$
(Moon). Writing $\boldsymbol r_1=\boldsymbol r-\boldsymbol r_E$ and
$\boldsymbol r_2=\boldsymbol r-\boldsymbol r_M$ with $r_i=\lVert\boldsymbol
r_i\rVert$, and denoting the ND maximum thrust acceleration $T$, exhaust velocity
$c=I_{sp}g_0$, throttle $s\in[0,1]$, and unit thrust direction $\boldsymbol
\alpha$, the fixed-time minimum-fuel problem is
\begin{equation}
\begin{aligned}
\min_{s(\cdot),\,\boldsymbol\alpha(\cdot)}\quad
  & J=\int_0^{t_f} s\,dt \\[2pt]
\text{s.t.}\quad
  & \dot{\boldsymbol r}=\boldsymbol v,\qquad
    \dot{\boldsymbol v}=\boldsymbol g(\boldsymbol r)+\boldsymbol h(\boldsymbol v)
      +\frac{sT}{m}\,\boldsymbol\alpha,\qquad
    \dot m=-\frac{T}{c}\,s,\\[2pt]
  & \boldsymbol g(\boldsymbol r)=\begin{bmatrix}x\\y\\0\end{bmatrix}
      -(1-\mu)\frac{\boldsymbol r_1}{r_1^{3}}
      -\mu\frac{\boldsymbol r_2}{r_2^{3}},\qquad
    \boldsymbol h(\boldsymbol v)=\begin{bmatrix}2\dot y\\-2\dot x\\0\end{bmatrix},\\[2pt]
  & s(t)\in[0,1],\qquad \lVert\boldsymbol\alpha(t)\rVert=1,\\[2pt]
  & \boldsymbol r(0),\boldsymbol v(0)=\text{(GTO departure)},\quad m(0)=1,\\[2pt]
  & \boldsymbol r(t_f),\boldsymbol v(t_f)=\text{(target arrival)},\quad
    m(t_f)\ \text{free},\quad t_f\ \text{fixed}.
\end{aligned}
\label{eq:ocp}
\end{equation}
Mass is normalized to $m(0)=1$, so the propellant burned is $m_p=1-m(t_f)=
\tfrac{T}{c}\int_0^{t_f}s\,dt$ and minimizing $J$ minimizes propellant directly.
The departure state is an equatorial GTO with perigee and apogee altitudes of
$350$~km and $35{,}786$~km and argument of perigee $-25^\circ$; the arrival state
is the tulip or ELFO rendezvous point of \cref{sec:constellation}. The two
targets differ only in these terminal conditions, so a single solver serves both.
The spacecraft modeled here is $15$~kg with a $25$~mN thruster at
$I_{sp}=2100$~s; because the transfer is governed by the thrust-to-mass ratio and
specific impulse, the result scales to heavier vehicles at fixed propellant mass
fraction. Solar radiation pressure, third-body perturbations beyond the two CR3BP
primaries, finite-attitude slew dynamics, and thruster duty-cycle constraints are
not modeled.

\subsection{Bang-bang structure and two sources of difficulty}
Minimum-fuel low thrust is the hard member of the minimum-time/minimum-energy/%
minimum-fuel family, and the reason is structural. The running cost $\int s\,dt$
is linear in the throttle, so Pontryagin's minimum principle drives $s$ onto its
bounds. With costates $\boldsymbol\lambda_r,\boldsymbol\lambda_v,\lambda_m$ and
Hamiltonian $H=\boldsymbol\lambda_r\!\cdot\!\boldsymbol v+\boldsymbol\lambda_v
\!\cdot\!\big(\boldsymbol g+\boldsymbol h+\tfrac{sT}{m}\boldsymbol\alpha\big)
-\lambda_m\tfrac{T}{c}s+s$, the thrust direction aligns with the primer vector
$\boldsymbol\alpha^\star=-\boldsymbol\lambda_v/\lVert\boldsymbol\lambda_v\rVert$
and $H$ is then affine in $s$. The optimal throttle is therefore bang-bang,
$s^\star=1$ where $S<0$ (burn) and $s^\star=0$ where $S>0$ (coast), with
switching function
\begin{equation}
S=1-\frac{T}{m}\lVert\boldsymbol\lambda_v\rVert-\frac{T}{c}\lambda_m ,
\label{eq:switch}
\end{equation}
and singular arcs occurring only on the measure-zero set $S\equiv0$. Over the
roughly 40-revolution GTO-to-cislunar spiral, $S$ changes sign near every perigee,
so the optimal control carries order-25 discontinuous switches.

Two independent difficulties compound at this optimum. The first is the
nonsmoothness just described: a control that is nowhere differentiable collapses
the convergence basin of indirect shooting and causes low-order collocation to
smear each switch, because a throttle jump between two mesh nodes registers as a
large defect that the optimizer rounds off. The second is dynamical and
independent of the objective: the $1/r_1^3$ gravity terms at the $\sim$40 perigee
passes, where $r_1\approx0.017$~ND, dominate the second derivatives and give the
multi-revolution trajectory a sensitivity of order $10^6$. Minimum-time faces
only the second difficulty and is comparatively easy, because maximum thrust is
optimal everywhere and its control has no switches; minimum-fuel faces both at
once. Each difficulty is removed by its own regularization below.

\subsection{Direct transcription}
The trajectory is transcribed by trapezoidal collocation on $N$ intervals: state
and control are carried at each node, the dynamics are enforced as trapezoidal
defect equalities between adjacent nodes, the endpoints are pinned as equality
constraints, and the objective is the trapezoidal quadrature of the running cost.
The unit-direction control is carried explicitly as $\boldsymbol\alpha$ with the
constraint $\lVert\boldsymbol\alpha\rVert=1$ alongside a separate throttle
$s\in[0,1]$, rather than as a single thrust vector; the vector form is degenerate
on coast arcs, where $s\to0$ leaves the direction undefined and wedges the solver.
The NLP is built symbolically in CasADi, which supplies the exact sparse
constraint Jacobian and Lagrangian Hessian by automatic differentiation, and is
solved with IPOPT. The exact Hessian is what resolves the many-switch coupling
that a limited-memory approximation cannot, but only after the dynamics are
regularized: on the raw problem the exact Hessian carries $1/r_1^5$ perigee terms
that destabilize the Newton step, and a solver using a damped limited-memory
Hessian is ironically more stable there. Regularize first, then the exact Hessian
becomes an asset rather than a liability.

\subsection{Sundman regularization of the perigee singularity}
The dynamics-side difficulty is removed by a Sundman change of independent
variable from time $t$ to a regularizing variable $\tau$,
\begin{equation}
\frac{dt}{d\tau}=\kappa(\boldsymbol r)=r_1^{\,p},\qquad p=1.5,
\label{eq:sundman}
\end{equation}
with every state equation multiplied by $\kappa$ and physical time carried as an
eighth state. This does two things at once. The near-perigee Hessian terms scale
as $r_1^{\,p-3}$ rather than $r_1^{-3}$, bounded for $p\ge3$ and already much
milder at $p=1.5$, which is what the exact-Hessian solve requires. And a uniform
mesh in $\tau$ is automatically dense in physical time wherever $\kappa=r_1^{\,p}$
is small, that is near perigee, exactly where the burns and switches concentrate,
so regularization and mesh refinement come from the same substitution. Carrying
$t$ as a state keeps the transfer time fixed while working in $\tau$: the terminal
condition $t(\tau_f)=t_f$ is imposed and the trajectory adjusts so that
$\int_0^{\tau_f}\kappa\,d\tau=t_f$. The regularized horizon $\tau_f$ is held fixed
rather than made a decision variable; a free $\tau_f$ multiplies every collocation
defect, producing a dense column in the constraint Jacobian that causes
catastrophic fill-in in the sparse linear solver and out-of-memory failure at
realistic mesh sizes.

The single-primary clock $\kappa=r_1^{\,p}$ regularizes only the Earth-perigee
singularity, which suffices for the tulip target but not for a rendezvous in the
Moon's vicinity, where the unresolved lunar gravity overflows the exact Hessian
at the near-Moon nodes. Replacing $\kappa$ with a two-primary minimum-distance
clock, $\kappa=(r_1^{-q}+(r_2/D)^{-q})^{-p/q}$ with $q\approx4$ and $D$ the ND
lunar sphere-of-influence radius ($\approx0.15$), recovers $r_1^{\,p}$ near Earth
and $(r_2/D)^{\,p}$ near the Moon and redistributes nodes into the capture arc.
The same regularization therefore serves both targets, differing only in whether
the second primary is active.

\subsection{Energy-to-fuel homotopy and warm starting}
The objective-side nonsmoothness is removed by continuation from a smooth
relaxation. The fuel cost is replaced by the Bertrand--\'Ep\'enoy family
\begin{equation}
J(\varepsilon)=\int_0^{t_f}\!\big[\,s-\varepsilon\,s(1-s)\,\big]\,dt,
\qquad\varepsilon\in[0,1].
\label{eq:homotopy}
\end{equation}
At $\varepsilon=0$ this is the fuel cost of \cref{eq:ocp}; at $\varepsilon=1$ it
is $\int s^2\,dt$, the minimum-energy cost, which is strictly convex in $s$, so
the optimal throttle is a smooth interior ramp rather than a bang-bang, the
costate dynamics are smooth, and the solver basin is large. The problem is solved
at $\varepsilon=1$ and then continued toward $\varepsilon=0$ in warm-started
steps; the ramp sharpens continuously into the bang-bang limit while the iterate
stays inside the basin. This ``start smooth, then sharpen'' continuation is the
single most effective element of the method.

Continuation requires a feasible starting trajectory, and how the seed is mapped
into the Sundman variable is decisive. Interpolating a time-mesh solution onto a
uniform $\tau$ mesh discards enough of the 40-revolution oscillatory structure to
leave an irreducible $\sim\!10^{-2}$ collocation defect, which the interior-point
solver reads as infeasibility and answers by remaining in its restoration phase,
exiting with a false ``locally infeasible.'' Mapping the seed onto its own nodes
instead, with each node's time carried to $\tau$ by trapezoidal quadrature of
\cref{eq:sundman}, leaves only a small trapezoid-rule mismatch, which the solver
closes to $\sim\!10^{-14}$ in normal mode. This no-resample seeding is what
converts the energy solve from perpetual restoration to convergence, and neither
regularization alone suffices: without the homotopy the Sundman solve thrashes
into restoration on the bang-bang objective, and without Sundman the homotopy
cannot drive the perigee defects below $\sim\!10^{-3}$.

\subsection{Optimality verification}
The direct transcription never forms the costates, but they can be recovered from
the converged NLP and checked against the minimum principle, because the
Karush--Kuhn--Tucker (KKT) multipliers of the defect constraints are the discrete
adjoint, up to a positive mesh-weight scaling and a global sign. Two of the
first-order conditions are scale-invariant and are checked directly. The thrust
direction of the direct solution matches the costate primer
$-\boldsymbol\lambda_v/\lVert\boldsymbol\lambda_v\rVert$ to a mean of
$0.06^\circ$ over all burn arcs, and the terminal mass costate, which
transversality requires to vanish under free final mass, is recovered at
$\sim\!10^{-7}$. That a solution obtained without ever forming costates
reproduces the primer direction to a hundredth of a degree is an independent
cross-check between the direct and indirect formulations. The remaining
first-order condition, the switching-function sign law of \cref{eq:switch}, is
scale-dependent and is verified separately by de-scaling the defect multipliers
to node costates; recovering the costates from the primal trajectory instead is
ill-posed here, because the forward adjoint map spans $\sim\!10^{11}$ over 40
revolutions and the fixed-endpoint rendezvous leaves $\boldsymbol\lambda_r$ and
$\boldsymbol\lambda_v$ without a boundary condition.

%==============================================================================
\section{Results}
\label{sec:results}
%==============================================================================
\emph{[A few sentences.]} Present the certified point solution, the two
$\Delta V$--time fronts, and the structural findings. The full-spiral tulip
min-fuel optimum is certified machine-tight and PMP-consistent; sweeping $t_f$
maps the propellant--trip-time trade for both targets and exposes a dense
local-minimum landscape and a hard near-min-time transition band.

\begin{itemize}
  \item \textbf{Certified tulip solution} ($1.15\times$): 25 switches, 99.4\%
        bang-bang, propellant 2.264~kg, $\Delta V=3.370$~km/s, max defect
        $2.4\times10^{-14}$, terminal error 0. Contrast with the 3-switch local
        optimum (2.950~kg) to show the global many-switch structure.
  \item \textbf{Tulip $\Delta V$--$t_f$ front}: certified band $1.12$--$1.25\times$
        (+$1.85\times$); campaign-best $\Delta V=2.434$~km/s at $1.65\times$
        (46~d, $\sim$45\% below min-time). Switch count grows/shrinks with
        $t_f$; dense-local-minima scatter and the lower envelope.
        \emph{[Outline note.]} ``Certified'' here $=$ PMP first-order via the
        KKT-dual costates (primer alignment + switching-law sign test); the
        low-$\Delta V$/fold points are well-reproduced \emph{feasible-envelope}
        bounds, deliberately not first-order-certified --- keep that distinction
        explicit wherever these numbers appear. \emph{[Also: the ``$\times$
        min-time'' factors and the ``\% below min-time'' figure are stated
        against $6.2907$~ND; recompute them against the target-consistent tulip
        anchor ($\sim$5.8~ND, \S2) --- physical $t_f$ and $\Delta V$ values are
        unchanged, only the min-time-relative labels shift.]}
  \item \textbf{Hard transition band} $1.01$--$1.11\times$: resists all methods;
        near-min-time conditioning wall that afflicts even the smooth energy
        problem---itself a structural finding.
  \item \textbf{ELFO $\Delta V$--$t_f$ front}: monotone $3.344$~km/s
        ($1.11\times$/31~d) $\to$ minimum $2.693$~km/s ($1.73\times$/48~d), then
        flat; optimum sits at a fold. 11/14 grid factors certified bang-bang.
        \emph{[Outline note.]} These ``$\times$'' were scaled by \emph{tulip's}
        $6.2907$~ND (the ELFO min-time was unsolved when the front was built);
        relabel against the now-certified ELFO min-time $6.0962$~ND. E.g.\ the
        $1.73\times$-tulip minimum $=1.79\times$ ELFO-min-time; $1.11\times$-tulip
        $=1.14\times$ ELFO. Same as tulip (\S2): physical $t_f$/$\Delta V$
        unchanged, only the labels.
  \item \textbf{Cross-target comparison / figures}: trajectory + throttle movie
        stills over the 40 revs; both fronts on one axis; PMP verification table.
\end{itemize}

%==============================================================================
\section{Future Work}
%==============================================================================
\emph{[A few sentences.]} Outline what remains: closing the transition band by
indirect/multiple-shooting from the min-time end, pinning exact fronts with
per-$t_f$ multi-start, and extending both the method and the mission scope.

\begin{itemize}
  \item Close the $1.01$--$1.11\times$ band: Sundman-domain multiple-shooting /
        switch-structure parameterization / saltation-aware (event-based)
        integration to remove the $1/\varepsilon$ smoothing layers.
  \item Accurate fronts: multi-start per $t_f$ to pin the lower envelope; fill
        the ELFO $1.65$--$2.0\times$ fold gaps.
  \item Indirect confirmation: full PMP TPBVP re-derivation as a
        belt-and-suspenders certificate.
  \item Mission extensions: additional targets, mass/thrust/$I_{sp}$ sweeps,
        variable-$t_f$ (free-time) objectives, station-keeping/re-phasing tie-in.
\end{itemize}

\end{document}
